\documentclass{article}

\usepackage{booktabs}
\usepackage{tabularx}
\usepackage{longtable}
\usepackage{array}
\usepackage{pdflscape}
\usepackage{hyperref}

\hypersetup{
    colorlinks=true,       % false: boxed links; true: colored links
    linkcolor=red,          % color of internal links (change box color with linkbordercolor)
    citecolor=green,        % color of links to bibliography
    filecolor=magenta,      % color of file links
    urlcolor=cyan           % color of external links
}

\title{Hazard Analysis\\\progname}

\author{\authname}

\date{}

\input{../Comments.text}
\input{../Common.text}

\begin{document}

\maketitle
\thispagestyle{empty}

~\newpage

\pagenumbering{roman}

\begin{table}[hp]
\caption{Revision History} \label{TblRevisionHistory}
\begin{tabularx}{\textwidth}{llX}
\toprule
\textbf{Date} & \textbf{Developer(s)} & \textbf{Change}\\
\midrule
2026-10-14 & Aaron Loh & First draft from team planning decisions\\
\bottomrule
\end{tabularx}
\end{table}

~\newpage

\tableofcontents

~\newpage

\pagenumbering{arabic}

\section{Introduction}

Natural Control lets blind users delegate tasks in a web browser and in Gmail and Google Calendar to an AI agent by voice. The user stays informed through spoken progress and results, and stays in control through clarification, approval, correction, and stopping.

\paragraph{Hazard definition.} A hazard is a property or condition of the system, together with a condition in its environment, that can lead to a loss (see section 2). Physical harm is not expected.

\paragraph{Purpose.} To find how each component could fail before we build it, and turn the safeguards into safety requirements for the SRS. A sighted user can glance at the screen to catch a mistake; a blind user often cannot, so detection has to be built into the system.

\section{Scope and Purpose of Hazard Analysis}

\paragraph{Included.} Voice interaction, task execution through the Hermes agent, browser tasks through Browser Use, and Gmail/Google Calendar actions, on a Mac in English, alongside VoiceOver.

\paragraph{Excluded.} Native Mac application control, general local-file tasks, phone calling, and saved workflows. These are deferred; each would add hazards to analyze if built.

\paragraph{Losses.} The hazards below can lead to these losses:

\begin{itemize}
  \item Unwanted action: something sent, submitted, purchased, created, or deleted that the user did not intend, including duplicates.
  \item Missed or failed task: not done while the user believes it was.
  \item Loss of control or awareness: the user cannot tell what the agent is doing, cannot stop it, or is told something false.
  \item Privacy or account misuse: private information exposed, or account access used beyond what the user intended.
\end{itemize}

\section{System Boundaries and Components}

Inputs are the user's voice and keyboard input. Outputs are actions on accounts and websites, and spoken feedback; no visual output is relied on. The system has seven components:

\begin{table}[h]
\centering
\caption{System components}\label{TblComponents}
\begin{tabularx}{\textwidth}{lX}
\toprule
\textbf{Component} & \textbf{Role}\\
\midrule
Accessible client (built) & Mic, stop key, and task history; runs alongside VoiceOver\\
Voice interface (LiveKit + our voice agent) & Hears, transcribes, interprets intent, and speaks responses\\
Task bridge (built) & Starts, tracks, stops, and resumes tasks in Hermes; keeps the task log\\
Hermes agent (reused) & Plans the task and calls tools, using an external AI model\\
Action gate (built) & Holds consequential actions for approval; verifies outcomes after acting\\
Browser Use (reused) & Reads and acts on websites\\
Account tools (built or adapted MCP server) & Gmail and Calendar access through OAuth\\
\bottomrule
\end{tabularx}
\end{table}

Websites, Google services (Gmail, Calendar), LiveKit, AI model and speech providers, network access, and VoiceOver. In: the user's voice and keyboard input. Out: actions on accounts and websites, and spoken feedback.

\todo{Add the component diagram.}

\section{Critical Assumptions}

\begin{itemize}
  \item Setup done correctly. A setup helper completed the assisted setup (accounts connected, permissions granted). Limits the analysis: setup errors are not analyzed as runtime hazards.
  \item One authorized speaker. Only the user gives commands; other voices are noise. Limits the analysis: we do not verify who is speaking.
  \item User's own accounts. Connected accounts belong to the user, who may act on them. Development and evaluation use test accounts.
\end{itemize}

We do not assume that speech recognition, the AI model, websites, or external services work correctly. Their failures appear in the FMEA.

\section{Failure Mode and Effect Analysis}

Each row is one failure mode of one component. Hazard IDs are H\textless{}component\textgreater{}-\textless{}letter\textgreater{}. How a failure is detected without sight is part of the recommended action, since a blind user often cannot notice a failure visually.

\begin{landscape}
\footnotesize
\setlength{\tabcolsep}{3pt}
\begin{longtable}{>{\raggedright\arraybackslash}p{0.9cm}>{\raggedright\arraybackslash}p{1.7cm}>{\raggedright\arraybackslash}p{2.9cm}>{\raggedright\arraybackslash}p{2.7cm}>{\raggedright\arraybackslash}p{3cm}>{\raggedright\arraybackslash}p{4.6cm}>{\raggedright\arraybackslash}p{1.6cm}}
\caption{Failure Mode and Effect Analysis}\label{TblFMEA}\\
\toprule
\textbf{Ref} & \textbf{Component} & \textbf{Failure mode} & \textbf{Effects} & \textbf{Causes} & \textbf{Recommended action} & \textbf{SR}\\
\midrule
\endfirsthead
\toprule
\textbf{Ref} & \textbf{Component} & \textbf{Failure mode} & \textbf{Effects} & \textbf{Causes} & \textbf{Recommended action} & \textbf{SR}\\
\midrule
\endhead
\bottomrule
\endfoot
H1-a & Accessible client & Stop control can't be reached without voice & User can't halt the agent if voice fails; loss of control & Shortcut clashes with VoiceOver; focus in another app; client frozen & Provide a stop shortcut that works whenever the client runs and doesn't clash with VoiceOver; test with VoiceOver on & FR.Sa.1\\
H1-b & Accessible client & Task history doesn't match what happened & User can't review or catch wrong actions afterward & History written from the agent's claims instead of recorded events; events lost during disconnects & Build history only from recorded actions and verified outcomes; label anything unverified & FR.Sa.2\\
H2-a & Voice interface & Misunderstands the user: wrong words, names, numbers, a cut-off sentence, or another voice & Wrong task or wrong details acted on & Background noise, accents, turn detection ending too early, TV or other people & Noise cancellation; read back key details (recipient, time, content summary) before every consequential action; allow spoken correction & FR.Sa.3\\
H2-b & Voice interface & A stop request isn't recognized & Task keeps running; loss of control & Stop classified as chatter; mistranscription; user talking over the system & Recognize stop phrases with a fixed rule before any model interprets them; keyboard stop as backup; acknowledge quickly & FR.Sa.4, NFR.Sa.1\\
H3-a & Task bridge & Stop is acknowledged but execution continues & Actions happen after the user said stop; false sense of safety & Hermes interrupts only at tool boundaries; an action is already in progress & Say ``stopping'' and ``stopped'' as separate messages; once a stop is requested, the gate refuses all further consequential actions for that task; report anything completed after the request & FR.Sa.5\\
H3-b & Task bridge & A consequential action is repeated after reconnecting & Duplicate email or event & Automatic retry after a dropped connection when the first attempt actually succeeded & Unique ID per approved action; never retry a consequential action automatically; pause new actions while disconnected; on reconnect, report pending items & FR.Sa.6, FR.Sa.7, NFR.Sa.2\\
H3-c & Task bridge & Task state or pending approvals lost & User never asked, or told the wrong status & Crash or restart of the bridge or Hermes & Save task state and pending approvals; after a restart, announce interrupted tasks and their last known status & FR.Sa.8\\
H4-a & Hermes agent & Does the wrong task or picks the wrong target & Wrong action prepared (wrong thread, wrong date) & Ambiguous request; model error & Ask a clarifying question when required details are missing or ambiguous; the approval read-back (H2-a) catches the rest & FR.Sa.9\\
H4-b & Hermes agent & Claims success that didn't happen & Missed task; user misinformed & Model hallucination; partial failure reported as success & Report completion only from verified outcomes; distinguish verified, partial, and unverified results & FR.Sa.10\\
H4-c & Hermes agent & Stalls or loops & Silence or repeated steps; user unsure what's happening; cost grows & Model loop; slow or broken website & Speak a status update at a regular interval while a task runs; after a step or time limit, pause and ask the user & FR.Sa.11, NFR.Sa.3\\
H4-d & Hermes agent & AI provider outage or timeout & Task halts partway & Provider down; network loss & Pause, tell the user what was and wasn't done, never claim success & FR.Sa.12\\
H5-a & Action gate & A consequential action slips through without approval & Unwanted action & Action not recognized as consequential (e.g. a browser click that submits a form); new tool not classified & Default-deny: anything not known to be read-only requires approval; if browser writes can't be intercepted reliably, make the browser read-only & FR.Sa.13\\
H5-b & Action gate & A stale approval is reused & Action runs with details the user never approved & Details changed after approval (time, recipient) & Bind each approval to its exact details; approvals are single-use; any change requires new approval & FR.Sa.14\\
H5-c & Action gate & An approval request expires or is withdrawn without a clear result & User unsure whether the action happened & Timeout; disconnect; request superseded & Tell the user explicitly that the action did not happen and what is still pending & FR.Sa.15\\
H5-d & Action gate & An irreversible delete is approved by mistake & Email or event permanently lost & Misheard approval; wrong item matched; user assumes the action can be undone & For deletes, read back identifying details (sender, subject, date, or event title and time) and say it cannot be undone; use recoverable deletion (trash) where the service supports it & FR.Sa.19\\
H6-a & Browser Use & Acts on the wrong page element & Wrong item chosen or wrong field filled & Poor labels; unusual layout; page changed after reading & Re-check the target against the approved details immediately before acting; verify the result afterward & FR.Sa.10, FR.Sa.14\\
H6-b & Browser Use & Follows instructions injected by a webpage & Unwanted action or private data sent elsewhere & Page or email text written to manipulate agents & Treat page and email content as data, never as commands; consequential actions must trace to the user's request; tell the user when content tries to give instructions & FR.Sa.16\\
H7-a & Account tools & Access is broader than needed & Any other failure can do more damage & Requesting broad OAuth permissions for convenience & Request the narrowest permissions each tool needs; keep tokens out of the repository and logs & FR.Sa.17\\
H7-b & Account tools & Outcome unknown after a timeout & Retrying risks a duplicate; not retrying risks a missed task & Network or API timeout & Check the real account state before any retry; if still unknown, tell the user and ask & FR.Sa.18\\
\end{longtable}
\end{landscape}

\section{Safety and Security Requirements}

These requirements come from the recommended actions in Table~\ref{TblFMEA} and are also part of the SRS. FR.Sa marks functional safety requirements and NFR.Sa non-functional ones. Timing values are named constants, set after POC measurements (see the SRS constants table).

\subsection{Functional Safety Requirements}

\begin{itemize}
  \item \textbf{FR.Sa.1:} The system shall provide a keyboard stop control that works whenever the client is running and does not conflict with VoiceOver commands. (H1-a)
  \item \textbf{FR.Sa.2:} The system shall build the task history only from recorded actions and verified outcomes, labelling unverified entries. (H1-b)
  \item \textbf{FR.Sa.3:} Before each consequential action, the system shall read back its key details and proceed only after explicit approval. (H2-a)
  \item \textbf{FR.Sa.4:} The system shall recognize stop requests with a fixed rule that does not depend on model interpretation. (H2-b)
  \item \textbf{FR.Sa.5:} After a stop request, the system shall refuse all further consequential actions for that task and announce stop acknowledged and stop confirmed as separate messages. (H3-a)
  \item \textbf{FR.Sa.6:} The system shall assign each approved consequential action a unique ID and shall not repeat an action with the same ID. (H3-b)
  \item \textbf{FR.Sa.7:} The system shall pause new actions while the user is disconnected, unless background continuation was authorized for that task. (H3-b)
  \item \textbf{FR.Sa.8:} The system shall save task state and pending approvals, and after a restart shall announce interrupted tasks with their last known status. (H3-c)
  \item \textbf{FR.Sa.9:} The system shall ask a clarifying question when details required for an action are missing or ambiguous. (H4-a)
  \item \textbf{FR.Sa.10:} The system shall report a task as complete only after verifying its outcome, and shall distinguish verified, partial, and unverified results. (H4-b, H6-a)
  \item \textbf{FR.Sa.11:} The system shall pause a task and ask the user after a configured step or time limit is reached. (H4-c)
  \item \textbf{FR.Sa.12:} If an external service fails mid-task, the system shall pause and tell the user what was and was not completed. (H4-d)
  \item \textbf{FR.Sa.13:} The system shall require approval for any action not classified as read-only. (H5-a)
  \item \textbf{FR.Sa.14:} The system shall bind each approval to the action's exact details, use it once, and require new approval if any detail changes. (H5-b, H6-a)
  \item \textbf{FR.Sa.15:} When an approval request expires or is withdrawn, the system shall tell the user that the action did not happen. (H5-c)
  \item \textbf{FR.Sa.16:} The system shall treat webpage and email content as data, and shall only perform consequential actions that trace to the user's request. (H6-b)
  \item \textbf{FR.Sa.17:} The system shall request only the account permissions each tool needs, and shall keep credentials out of source control and logs. (H7-a)
  \item \textbf{FR.Sa.18:} The system shall check the real account state before retrying an action whose outcome is unknown, and shall ask the user if it remains unknown. (H7-b)
  \item \textbf{FR.Sa.19:} For delete actions, the system shall read back the item's identifying details, state that the action cannot be undone, and use recoverable deletion (such as moving to trash) where the service supports it. (H5-d)
\end{itemize}

\subsection{Non-Functional Safety Requirements}

\begin{itemize}
  \item \textbf{NFR.Sa.1:} The system shall audibly acknowledge a stop request within STOP\_ACK\_TIME of the end of the utterance or key press. (H2-b)
  \item \textbf{NFR.Sa.2:} The system shall notify the user of a lost connection within CONNECTION\_NOTICE\_TIME, by sound if speech is unavailable. (H3-b)
  \item \textbf{NFR.Sa.3:} While a task is running, the system shall give a spoken status update after IDLE\_PROMPT\_TIME without spoken progress. (H4-c)
\end{itemize}

\section{Roadmap}

\paragraph{POC (November).} FR.Sa.3, 4, 5, 13, 14; NFR.Sa.1. The safeguards the demo is meant to prove. FR.Sa.13 includes testing whether browser actions can be intercepted.

\paragraph{Before Rev 0 / participant evaluation.} FR.Sa.1, 2, 6, 7, 8, 9, 10, 12, 15, 16, 17, 18, 19; NFR.Sa.2. These matter as soon as real participants use real-looking accounts.

\paragraph{Later.} FR.Sa.11; NFR.Sa.3. Basic status updates exist earlier, but tuned step/time limits and update frequency need usability evidence. (Team to confirm.)

\newpage{}

\section{Appendix: Generative AI Use Disclosure}

\subsection{AI Tools Used}

Aaron Loh: Claude Code with Claude Opus 5.5.

\todo{Other members: add the tools you used.}

\subsection{How and Where Used}

Aaron used Claude to work through this analysis. It asked questions about scope, losses, components, and failure modes, summarized the Hermes, LiveKit, and Browser Use documentation, and translated the decisions into written drafts of each section. Aaron made and reviewed every decision.

\subsection{Verification of AI-Generated Content}

This is a first draft. Before submission the team will review every hazard and requirement, check the cited documentation, rewrite sections in our own words, and keep this disclosure up to date. The team is responsible for the submitted document.

\newpage{}

\section{Appendix: Reflection}

\input{../Reflection.text}

\begin{enumerate}

    \item What went well while writing this deliverable? 
    \item What pain points did you experience during this deliverable, and how
    did you resolve them?
    \item Which of your listed risks had your team thought of before this
    deliverable, and which did you think of while doing this deliverable? For
    the latter ones (ones you thought of while doing the Hazard Analysis), how
    did they come about?
    \item Other than the risk of physical harm (some projects may not have any
    appreciable risks of this form), list at least 2 other types of risk in
    software products. Why are they important to consider?
\end{enumerate}

\end{document}